% Part IV -- Numerical Data, Tables, and Visualization
% Chapter 13: Visualizing data with Matplotlib

\h{Visualizing Data with Matplotlib}
\label{ch:13}

A chart should make one question easier to answer. Matplotlib supplies many
plot types, but a beginner needs only a small, reliable set: line plots for
ordered change, bar charts for category comparisons, scatter plots for
relationships, and histograms for distributions. This chapter builds those
charts through one consistent object-based pattern.

\begin{NoteBox}
\textbf{Prerequisites.} You should understand lists, NumPy arrays, pandas
DataFrames, column selection, and basic numerical summaries.
\end{NoteBox}

\begin{tcolorbox}[title=Learning outcomes,colframe=orange,colback=white,arc=2mm]
After completing this chapter, you should be able to:
\begin{itemize}
    \item create a Matplotlib Figure and Axes with \c{plt.subplots()};
    \item choose a line, bar, scatter, or histogram from the question and data;
    \item add a clear title, axis labels, units, legend, and restrained grid;
    \item use a consistent palette without depending on color alone;
    \item plot selected NumPy or pandas data;
    \item arrange two related charts as subplots; and
    \item save a figure with an appropriate file name and bounding box.
\end{itemize}
\end{tcolorbox}

\hh{Install once and import when needed}

Matplotlib is an external library. Install it once in the active Python
environment by running this command in a terminal.

\begin{verbatim}
python -m pip install matplotlib
\end{verbatim}

Use the same Python command that runs your scripts. Replace \texttt{python}
with \texttt{py} or \texttt{python3} if that is the working command on your
computer.

Most introductory plotting uses the \c{pyplot} interface with the conventional
name \c{plt}.

\begin{lstlisting}[
    language=python,
    style=block,
    caption={Create and close an empty figure},
    label={c13_import_matplotlib},
    captionpos=b
]
import matplotlib.pyplot as plt

fig, ax = plt.subplots()
print(type(fig).__name__)
print(type(ax).__name__)
plt.close(fig)
\end{lstlisting}

Closing a finished figure releases its resources. This matters when a program
creates many figures in one run.

\begin{SyntaxBox}{Create the container before drawing the chart}
The statement \c{fig, ax = plt.subplots()} calls one function and unpacks its
two returned objects. Use \c{ax} for plotting and labels, use \c{fig} for the
complete output such as \c{fig.savefig(...)}, and use \c{plt.show()} only when
you want an interactive display window.
\end{SyntaxBox}

\hh{Understand Figure and Axes}

A \c{Figure} is the complete output area. An \c{Axes} is one plotting region
inside that Figure. Most visible chart elements are added through an Axes.

\begin{lstlisting}[
    language=python,
    style=block,
    caption={Draw a first line plot},
    label={c13_first_plot},
    captionpos=b
]
import matplotlib.pyplot as plt

days = [1, 2, 3, 4, 5]
temperatures_c = [19.0, 20.5, 22.0, 21.0, 23.5]

fig, ax = plt.subplots(layout="constrained")
ax.plot(days, temperatures_c)

plt.show()
plt.close(fig)
\end{lstlisting}

The call to \c{plt.subplots()} creates both objects. The call to \c{ax.plot()}
adds data to the Axes. In a standalone script, \c{plt.show()} displays the
Figure in a window when the active backend supports one.

\begin{ExplainBox}{A chart maps variables to visual channels}
A visualization is a rule that maps data values to position, length, angle,
area, color, marker shape, or another visible property. Position on a common
scale usually supports more precise comparison than area or color. Use color
or marker shape to distinguish categories, not to replace a labeled numerical
axis. Before writing plotting syntax, state which variable controls each visual
channel and what comparison the reader should be able to make.
\end{ExplainBox}

The first plot maps \c{days} to horizontal position and
\c{temperatures_c} to vertical position. Connecting the positions with a line
adds a claim that order and continuity between adjacent days matter. If the
x-values were unrelated room categories, bars or points would express the
comparison more honestly than a connecting line.

\hh{Add the information a reader needs}

An unlabeled line has little meaning. Name the subject, both variables, and
their units.

\begin{lstlisting}[
    language=python,
    style=block,
    caption={Label a line plot clearly},
    label={c13_labeled_line},
    captionpos=b
]
import matplotlib.pyplot as plt

days = [1, 2, 3, 4, 5]
temperatures_c = [19.0, 20.5, 22.0, 21.0, 23.5]

fig, ax = plt.subplots(layout="constrained")
ax.plot(days, temperatures_c, color="tab:blue", marker="o")
ax.set_title("Indoor temperature over five days")
ax.set_xlabel("Day")
ax.set_ylabel("Temperature (degrees C)")
ax.grid(axis="y", alpha=0.25)

plt.show()
plt.close(fig)
\end{lstlisting}

\bookfigure[0.78\textwidth]{Figures/Generated/c13_labeled_line.pdf}
{Rendered result of Listing~\ref{c13_labeled_line}: labels and units make the plotted values interpretable.}
{fig:c13-labeled-line-result}

The grid is light and appears only on the value axis. It supports reading
values without competing with the data.

\begin{ExplainBox}{Read the result as well as the syntax}
The markers identify five measured days. The line shows their ordered pattern:
temperature rises overall, with a small decrease on day 4. That interpretation
depends on the title and axis labels; the code alone is not the final
communication product.
\end{ExplainBox}

\hh{Choose a plot from the question}

Start with the relationship you want to show.

\begin{center}
\begin{tabular}{lll}
\textbf{Question} & \textbf{Typical data} & \textbf{Plot} \\
How does a value change in order? & Time or distance sequence & Line \\
How do categories compare? & Names and values & Bar \\
Do two numerical variables relate? & Paired measurements & Scatter \\
How are values distributed? & One numerical collection & Histogram \\
\end{tabular}
\end{center}

Do not select a chart because it looks impressive. Select it because its visual
structure matches the question.

\begin{TryBox}{Match questions to visual structures}
Write one question about change over time, one about category differences, one
about a relationship between two measurements, and one about the distribution
of a single measurement. Name the plot type for each question before creating
any code.
\end{TryBox}

\hh{Visualization showcase: one library, many questions}
\label{sec:c13-showcase}

The course introduction shows that visualization extends well beyond four
introductory plot types. Figure~\ref{fig:c13-visualization-showcase} presents a
curated set of cases that recur in design analysis, building performance,
optimization, and computational geometry. Treat the gallery as a question
map, not as a list of effects to add indiscriminately.

\bookfigure[0.98\textwidth]{Figures/Generated/c13_visualization_showcase.pdf}
{A visualization showcase connecting six analytical questions to six visual structures.}
{fig:c13-visualization-showcase}

\begin{itemize}
    \item \textbf{Trend and uncertainty}: the line shows the central estimate;
          the translucent band shows a plausible range.
    \item \textbf{Annotated heatmap}: position identifies a row and column,
          color encodes magnitude, and text preserves exact values.
    \item \textbf{Distribution shapes}: each violin summarizes spread,
          concentration, and median rather than showing only an average.
    \item \textbf{Composition by floor}: stacked areas show how several
          components contribute to a changing total.
    \item \textbf{Optimization landscape}: contours reveal regions of similar
          objective values and the location of a candidate solution.
    \item \textbf{Three-dimensional response}: a surface represents a result
          that depends on two continuous inputs; it should be used only when
          the third dimension clarifies the relationship.
\end{itemize}

\begin{SyntaxBox}{The Figure/Axes pattern scales to advanced views}
The plotting object remains an Axes even when the method changes. Use
\c{ax.fill_between()} for a band, \c{ax.imshow()} for a matrix,
\c{ax.violinplot()} for distributions, \c{ax.stackplot()} for changing
composition, and \c{ax.contourf()} for filled contours. A 3D Axes is created
with \c{fig.add_subplot(..., projection="3d")}; its
\c{plot_surface()} method receives two coordinate grids and one result grid.
\end{SyntaxBox}

\hh{Use a line plot for ordered change}

A line connects observations in sequence. That connection implies meaningful
order, such as time, distance, or iteration.

\begin{lstlisting}[
    language=python,
    style=block,
    caption={Show ordered change with a line},
    label={c13_line_plot},
    captionpos=b
]
import matplotlib.pyplot as plt

hours = [8, 10, 12, 14, 16, 18]
daylight_percent = [18, 35, 52, 61, 43, 20]

fig, ax = plt.subplots(layout="constrained")
ax.plot(
    hours,
    daylight_percent,
    color="tab:blue",
    marker="o",
    linewidth=2,
    label="Measured",
)
ax.set(
    title="Daylight through the day",
    xlabel="Hour",
    ylabel="Daylight level (percent)",
)
ax.legend()
ax.grid(axis="y", alpha=0.25)

plt.show()
plt.close(fig)
\end{lstlisting}

\bookfigure[0.78\textwidth]{Figures/Generated/c13_line_plot.pdf}
{Rendered line plot: the measured daylight level rises toward 14:00 and then falls.}
{fig:c13-line-plot-result}

Markers reveal the actual observations. The connecting line helps the eye
follow their order but does not create new measurements between them.

\hh{Compare two ordered series}

Give each series both a label and a visual pattern. Color alone should not be
the only way to tell them apart.

\begin{lstlisting}[
    language=python,
    style=block,
    caption={Compare two series with redundant visual cues},
    label={c13_two_lines},
    captionpos=b
]
import matplotlib.pyplot as plt

hours = [8, 10, 12, 14, 16, 18]
north_c = [19.0, 20.0, 21.0, 21.5, 21.0, 20.0]
south_c = [19.5, 21.0, 23.0, 24.0, 22.5, 21.0]

fig, ax = plt.subplots(layout="constrained")
ax.plot(
    hours,
    north_c,
    color="tab:blue",
    marker="o",
    linestyle="-",
    label="North room",
)
ax.plot(
    hours,
    south_c,
    color="tab:orange",
    marker="s",
    linestyle="--",
    label="South room",
)
ax.set(
    title="Room temperature comparison",
    xlabel="Hour",
    ylabel="Temperature (degrees C)",
)
ax.legend()
ax.grid(axis="y", alpha=0.25)

plt.show()
plt.close(fig)
\end{lstlisting}

\bookfigure[0.78\textwidth]{Figures/Generated/c13_two_lines.pdf}
{Rendered comparison: the south room becomes warmer and varies more during the day.}
{fig:c13-two-lines-result}

The names in the legend carry the meaning. Marker shape and line style keep the
series distinguishable when color is hard to perceive or the page is printed
in grayscale.

\hh{Use a bar chart for categories}

Bars compare values attached to separate categories. Their lengths share a
common baseline, which makes differences easy to judge.

\begin{lstlisting}[
    language=python,
    style=block,
    caption={Compare category values with bars},
    label={c13_bar_chart},
    captionpos=b
]
import matplotlib.pyplot as plt

room_names = ["Lobby", "Office", "Cafe", "Studio"]
areas_m2 = [24.0, 18.0, 31.0, 26.0]

fig, ax = plt.subplots(layout="constrained")
bars = ax.bar(room_names, areas_m2, color="tab:blue")
ax.bar_label(bars, fmt="%.1f", padding=3)
ax.set(
    title="Area by room",
    xlabel="Room",
    ylabel="Area (m2)",
)
ax.set_ylim(0, 36)
ax.grid(axis="y", alpha=0.25)
ax.set_axisbelow(True)

plt.show()
plt.close(fig)
\end{lstlisting}

\bookfigure[0.74\textwidth]{Figures/Generated/c13_bar_chart.pdf}
{Rendered bar chart: a shared zero baseline supports an honest comparison of room areas.}
{fig:c13-bar-chart-result}

A zero baseline keeps bar lengths honest. Direct value labels are useful when
the exact values matter and the chart has only a few bars.

\hh{Use horizontal bars for long labels}

Horizontal bars give long category names more room.

\begin{lstlisting}[
    language=python,
    style=block,
    caption={Use horizontal bars for readable category labels},
    label={c13_horizontal_bar},
    captionpos=b
]
import matplotlib.pyplot as plt

materials = [
    "Reinforced concrete",
    "Cross-laminated timber",
    "Structural steel",
]
quantities_m3 = [42.0, 28.0, 16.0]

fig, ax = plt.subplots(layout="constrained")
ax.barh(materials, quantities_m3, color="tab:orange")
ax.set(
    title="Material quantities",
    xlabel="Quantity (m3)",
    ylabel="Material",
)
ax.grid(axis="x", alpha=0.25)
ax.set_axisbelow(True)

plt.show()
plt.close(fig)
\end{lstlisting}

\bookfigure[0.78\textwidth]{Figures/Generated/c13_horizontal_bar.pdf}
{Rendered horizontal bars leave space for long, readable material names.}
{fig:c13-horizontal-bar-result}

Long labels remain horizontal, so the reader does not need to rotate the page
or decode abbreviations.

\hh{Use a scatter plot for paired measurements}

A scatter plot represents each observation as one point with an x-value and a
y-value. It can reveal direction, spread, clusters, and unusual observations.

\begin{lstlisting}[
    language=python,
    style=block,
    caption={Show a relationship between two numerical variables},
    label={c13_scatter_plot},
    captionpos=b
]
import matplotlib.pyplot as plt

window_percent = [18, 22, 28, 35, 41, 48, 55]
daylight_percent = [24, 29, 37, 45, 53, 58, 66]

fig, ax = plt.subplots(layout="constrained")
ax.scatter(
    window_percent,
    daylight_percent,
    color="tab:blue",
    edgecolor="black",
    alpha=0.8,
)
ax.set(
    title="Window area and daylight",
    xlabel="Window area (percent of wall)",
    ylabel="Daylight level (percent)",
)
ax.grid(alpha=0.2)

plt.show()
plt.close(fig)
\end{lstlisting}

\bookfigure[0.76\textwidth]{Figures/Generated/c13_scatter_plot.pdf}
{Rendered scatter plot: each point preserves a pair of window-area and daylight measurements.}
{fig:c13-scatter-result}

The pattern suggests an association, not proof that one variable causes the
other. Other variables may explain part of the relationship.

\begin{NoteBox}
\textbf{Stop and predict.} What would a point far above the other points mean
in this example? Describe it using both variables before explaining a possible
cause.
\end{NoteBox}

\hh{Use a histogram for a distribution}

A histogram groups a numerical collection into intervals called bins. It shows
where values concentrate and how broadly they spread.

\begin{lstlisting}[
    language=python,
    style=block,
    caption={Group numerical values into histogram bins},
    label={c13_histogram},
    captionpos=b
]
import matplotlib.pyplot as plt

room_areas_m2 = [
    15, 18, 18, 20, 21, 22, 22, 23, 24, 24,
    25, 26, 27, 28, 30, 31, 33, 35, 38, 42,
]

fig, ax = plt.subplots(layout="constrained")
ax.hist(
    room_areas_m2,
    bins=[10, 20, 30, 40, 50],
    color="tab:blue",
    edgecolor="black",
)
ax.set(
    title="Distribution of room areas",
    xlabel="Area (m2)",
    ylabel="Number of rooms",
)
ax.grid(axis="y", alpha=0.25)
ax.set_axisbelow(True)

plt.show()
plt.close(fig)
\end{lstlisting}

\bookfigure[0.76\textwidth]{Figures/Generated/c13_histogram.pdf}
{Rendered histogram: bin heights show how many rooms fall inside each area interval.}
{fig:c13-histogram-result}

Different bins can change the visible pattern. State or inspect bin boundaries
when the conclusion depends on them.

\hh{Use color consistently}

This book uses orange, black, and white for instructional page elements. Chart
colors serve a different purpose: they distinguish data while remaining
restrained.

\begin{center}
\begin{tabular}{lll}
\textbf{Matplotlib color} & \textbf{Typical role} & \textbf{Additional cue} \\
\c{tab:blue} & Primary series & Solid line or circle marker \\
\c{tab:orange} & Comparison series & Dashed line or square marker \\
\c{tab:green} & Optional third series & Dotted line or triangle marker \\
\end{tabular}
\end{center}

Keep the same meaning for a color within one figure. Use text labels, a legend,
marker shapes, or line styles as well. Avoid adding colors that carry no
information.

\hh{Plot data from a pandas DataFrame}

Select columns by name, then pass those Series to an Axes method. This keeps
data preparation separate from chart formatting.

\begin{lstlisting}[
    language=python,
    style=block,
    caption={Prepare pandas data before plotting},
    label={c13_plot_dataframe},
    captionpos=b
]
import matplotlib.pyplot as plt
import pandas as pd

rooms = pd.DataFrame({
    "room": ["Lobby", "Office", "Cafe", "Studio"],
    "area_m2": [24.0, 18.0, 31.0, 26.0],
})
ordered = rooms.sort_values("area_m2")

fig, ax = plt.subplots(layout="constrained")
ax.barh(
    ordered["room"],
    ordered["area_m2"],
    color="tab:blue",
)
ax.set(
    title="Room areas from a DataFrame",
    xlabel="Area (m2)",
    ylabel="Room",
)
ax.grid(axis="x", alpha=0.25)
ax.set_axisbelow(True)

plt.show()
plt.close(fig)
\end{lstlisting}

\bookfigure[0.76\textwidth]{Figures/Generated/c13_dataframe_plot.pdf}
{Rendered pandas-based plot: sorting before plotting makes the category ranking visible.}
{fig:c13-dataframe-plot-result}

Sorting is part of the data-preparation stage. The plotting stage receives the
rows in the intended visual order.

\hh{Create two related subplots}

Subplots help readers compare related views that should share one Figure. Keep
the number small and give each Axes a distinct question.

\begin{lstlisting}[
    language=python,
    style=block,
    caption={Place two related questions in one Figure},
    label={c13_two_subplots},
    captionpos=b
]
import matplotlib.pyplot as plt
import pandas as pd

data = pd.DataFrame({
    "month": ["Jan", "Feb", "Mar", "Apr", "May", "Jun"],
    "energy_kwh": [820, 760, 690, 610, 540, 500],
    "temperature_c": [2, 4, 8, 13, 18, 22],
})

fig, axes = plt.subplots(1, 2, figsize=(10, 4), layout="constrained")

axes[0].plot(
    data["month"],
    data["energy_kwh"],
    color="tab:blue",
    marker="o",
)
axes[0].set(
    title="Monthly energy",
    xlabel="Month",
    ylabel="Energy (kWh)",
)

axes[1].scatter(
    data["temperature_c"],
    data["energy_kwh"],
    color="tab:orange",
    marker="s",
    edgecolor="black",
)
axes[1].set(
    title="Energy and temperature",
    xlabel="Outdoor temperature (degrees C)",
    ylabel="Energy (kWh)",
)

for ax in axes:
    ax.grid(alpha=0.2)

plt.show()
plt.close(fig)
\end{lstlisting}

\bookfigure[0.94\textwidth]{Figures/Generated/c13_two_subplots.pdf}
{Rendered two-panel Figure: the first Axes preserves time order while the second emphasizes a numerical relationship.}
{fig:c13-two-subplots-result}

The first chart preserves time order. The second removes that ordering to focus
on the relationship between two numerical variables.

\hh{Applied example: a building-performance dashboard}
\label{sec:c13-dashboard}

A dashboard should coordinate several questions without repeating the same
view. This example combines monthly composition, temperature context, a weekly
occupancy matrix, and total monthly energy. The script is intentionally longer
than earlier examples so you can practice reading a realistic program in
stages.

A runnable copy is provided in
\c{examples/c13_building_performance_dashboard.py} so you can modify the data
and panels without retyping the listing.

\begin{lstlisting}[
    language=python,
    style=block,
    caption={Build and export a four-panel building-performance dashboard},
    label={c13_building_dashboard},
    captionpos=b
]
from pathlib import Path

import matplotlib.pyplot as plt
import numpy as np

# 1. Prepare repeatable example data.
rng = np.random.default_rng(12)
months = np.arange(1, 13)
heating = np.array([52, 46, 35, 22, 12, 5, 3, 4, 10, 24, 38, 49])
cooling = np.array([2, 3, 5, 9, 18, 31, 38, 36, 24, 12, 5, 2])
lighting = np.array([18, 17, 16, 15, 14, 13, 13, 13, 14, 15, 17, 18])
equipment = np.full(12, 21)
outdoor_c = np.array([-2, 0, 5, 11, 17, 22, 25, 24, 20, 13, 7, 1])

occupancy = np.clip(rng.normal(38, 22, size=(7, 24)), 0, 100)
occupancy[:, :7] *= 0.15
occupancy[:, 19:] *= 0.25
occupancy[5:, :] *= 0.35

# 2. Create a 2-by-2 array of Axes.
fig, axes = plt.subplots(2, 2, figsize=(11, 7), layout="constrained")
fig.suptitle("Building performance dashboard")

# 3. Show monthly composition.
axes[0, 0].stackplot(
    months,
    heating,
    cooling,
    lighting,
    equipment,
    labels=["Heating", "Cooling", "Lighting", "Equipment"],
    colors=["tab:blue", "tab:orange", "0.60", "#C6A15B"],
    alpha=0.88,
)
axes[0, 0].set(
    title="Monthly energy by end use",
    xlabel="Month",
    ylabel="Energy (MWh)",
)
axes[0, 0].legend(fontsize=8, ncol=2)

# 4. Add temperature context and a reference band.
axes[0, 1].plot(months, outdoor_c, color="tab:blue", marker="o")
axes[0, 1].axhspan(18, 24, color="tab:green", alpha=0.14)
axes[0, 1].set(
    title="Outdoor temperature context",
    xlabel="Month",
    ylabel="Temperature (degrees C)",
)

# 5. Convert the 7-by-24 matrix into a heatmap.
heatmap = axes[1, 0].imshow(
    occupancy,
    cmap="YlOrRd",
    aspect="auto",
    vmin=0,
    vmax=100,
)
axes[1, 0].set(
    title="Weekly occupancy pattern",
    xlabel="Hour",
    ylabel="Day",
    yticks=range(7),
    yticklabels=["Mon", "Tue", "Wed", "Thu", "Fri", "Sat", "Sun"],
)
fig.colorbar(
    heatmap,
    ax=axes[1, 0],
    orientation="horizontal",
    fraction=0.08,
    pad=0.14,
    label="Occupancy (percent)",
)

# 6. Calculate and compare totals.
totals = heating + cooling + lighting + equipment
axes[1, 1].bar(months, totals, color="tab:blue")
axes[1, 1].axhline(
    totals.mean(),
    color="tab:orange",
    linestyle="--",
    label=f"Mean {totals.mean():.1f}",
)
axes[1, 1].set(
    title="Total monthly energy",
    xlabel="Month",
    ylabel="Energy (MWh)",
)
axes[1, 1].legend()

# 7. Apply shared finishing steps and export.
for ax in axes.flat:
    ax.grid(alpha=0.18)

output_folder = Path("study_figures")
output_folder.mkdir(exist_ok=True)
fig.savefig(
    output_folder / "building_performance_dashboard.png",
    dpi=200,
    bbox_inches="tight",
)
plt.show()
plt.close(fig)
\end{lstlisting}

\bookfigure[0.96\textwidth]{Figures/Generated/c13_building_dashboard.pdf}
{Rendered dashboard from Listing~\ref{c13_building_dashboard}. Each panel answers a different but related building-performance question.}
{fig:c13-building-dashboard}

\begin{SyntaxBox}{Read a long plotting script by stages}
\begin{enumerate}
    \item Imports identify the external tools and the path API used for export.
    \item Named arrays separate data preparation from drawing commands.
    \item \c{axes[row, column]} selects one panel from the 2-by-2 Axes array.
    \item Each plotting call returns or creates an artist. The heatmap artist is
          stored because \c{fig.colorbar()} needs it to build a matching scale.
    \item \c{axes.flat} provides one sequence for shared formatting instead of
          repeating the same grid instruction four times.
    \item Export occurs only after every panel is complete; the Figure, not one
          Axes, calls \c{savefig()}.
\end{enumerate}
\end{SyntaxBox}

Read the dashboard as an argument. Heating and cooling exchange seasonal
importance, the temperature panel supplies environmental context, the heatmap
shows when the building is occupied, and the total-energy bars make the yearly
cycle easy to compare. A dashboard is successful when these views support one
another without requiring the reader to decode decorative complexity.

\hh{Annotate one important observation}

An annotation should direct attention to a value relevant to the chart's
question. Too many annotations become clutter.

\begin{lstlisting}[
    language=python,
    style=block,
    caption={Annotate one observation without crowding the plot},
    label={c13_annotation},
    captionpos=b
]
import matplotlib.pyplot as plt

months = ["Jan", "Feb", "Mar", "Apr", "May", "Jun"]
energy_kwh = [820, 760, 690, 610, 540, 500]

fig, ax = plt.subplots(layout="constrained")
ax.plot(months, energy_kwh, color="tab:blue", marker="o")
ax.annotate(
    "Lowest value",
    xy=("Jun", 500),
    xytext=("Apr", 700),
    arrowprops={"arrowstyle": "->", "color": "black"},
)
ax.set(
    title="Monthly energy use",
    xlabel="Month",
    ylabel="Energy (kWh)",
)
ax.grid(axis="y", alpha=0.25)

plt.show()
plt.close(fig)
\end{lstlisting}

\bookfigure[0.78\textwidth]{Figures/Generated/c13_annotation.pdf}
{Rendered annotation: one concise callout directs attention to the lowest energy value.}
{fig:c13-annotation-result}

The annotation text states the interpretation. Its arrow identifies the
specific observation.

\hh{Save a figure}

Call \c{savefig()} on the Figure. The file extension selects the output format.
PNG is convenient for screens; PDF and SVG preserve vector lines and text.

\begin{lstlisting}[
    language=python,
    style=block,
    caption={Save a labeled figure to a folder},
    label={c13_save_figure},
    captionpos=b
]
from pathlib import Path
import matplotlib.pyplot as plt

output_folder = Path("study_figures")
output_folder.mkdir(exist_ok=True)
output_path = output_folder / "area_by_room.png"

fig, ax = plt.subplots(layout="constrained")
ax.bar(
    ["Lobby", "Office", "Cafe"],
    [24.0, 18.0, 31.0],
    color="tab:blue",
)
ax.set(
    title="Area by room",
    xlabel="Room",
    ylabel="Area (m2)",
)

fig.savefig(output_path, dpi=200, bbox_inches="tight")
plt.close(fig)

print(f"Saved: {output_path}")
\end{lstlisting}

The \c{bbox_inches="tight"} option trims unnecessary outer whitespace while
keeping labels in the output. Always open the saved file and inspect it before
sharing or printing it.

\hh{Use a repeatable chart workflow}

A simple order prevents formatting from hiding a data problem.

\begin{enumerate}
    \item State the question in one sentence.
    \item Inspect and prepare the data.
    \item Choose the plot type from the variable relationship.
    \item Create the Figure and Axes.
    \item Plot the data.
    \item Add title, labels, units, and a legend when needed.
    \item Reduce clutter and check accessibility.
    \item Save or show the Figure, then inspect the result.
\end{enumerate}

This order is short enough to use every time.

\hh{Common mistakes}

\begin{itemize}
    \item \textbf{Choosing by appearance}: match the plot type to the question
    and variable types.
    \item \textbf{Omitting units}: include units in axis labels whenever values
    represent measurements.
    \item \textbf{Using a line for unordered categories}: lines imply a
    meaningful sequence; use bars for separate categories.
    \item \textbf{Starting a bar axis above zero}: a truncated baseline can
    exaggerate differences in bar length.
    \item \textbf{Depending on color alone}: add labels, markers, or line
    styles.
    \item \textbf{Adding heavy decoration}: use light grids and remove elements
    that do not help answer the question.
    \item \textbf{Forgetting to inspect the saved file}: rendering and cropping
    can differ from the interactive window.
\end{itemize}

\hh{Practice ladder}

\hhh{Level 0 -- Read and predict}

Which plot type is produced? Which variable appears on each axis?

\begin{lstlisting}[
    language=python,
    style=block,
    caption={Identify a plot type and its variables},
    label={c13_practice_level0},
    captionpos=b
]
import matplotlib.pyplot as plt

widths_m = [3.0, 4.0, 5.0]
areas_m2 = [18.0, 24.0, 32.0]

fig, ax = plt.subplots()
ax.scatter(widths_m, areas_m2)
plt.close(fig)
\end{lstlisting}

\hhh{Level 1 -- Modify}

Add a title, both axis labels with units, and value labels above the bars.

\begin{lstlisting}[
    language=python,
    style=block,
    caption={Complete the communication around a bar chart},
    label={c13_practice_level1},
    captionpos=b
]
import matplotlib.pyplot as plt

rooms = ["A", "B", "C"]
areas_m2 = [18.0, 24.0, 31.0]

fig, ax = plt.subplots(layout="constrained")
ax.bar(rooms, areas_m2, color="tab:blue")
plt.show()
plt.close(fig)
\end{lstlisting}

\hhh{Level 2 -- Complete}

Create a line plot for six ordered temperature readings. Use circle markers,
a solid \c{tab:blue} line, a light horizontal grid, and labels with units.

\hhh{Level 3 -- Apply}

Create a two-panel Figure from a pandas DataFrame. The left Axes should compare
category values with bars. The right Axes should show the distribution of the
same numerical values with a histogram.

\textbf{Hints}
\begin{itemize}
    \item \textbf{Hint 1}: Create the Figure with \c{plt.subplots(1, 2)}.
    \item \textbf{Hint 2}: Select DataFrame columns by name.
    \item \textbf{Hint 3}: Give each Axes its own title and labels.
\end{itemize}

\hhh{Level 4 -- Challenge}

Write \c{plot_two_series(x, first, second, labels)}. It should return a Figure
with two labeled lines. Use different colors, markers, and line styles; include
axis labels, units, a legend, and a light grid. Test the function and close the
returned Figure.

\hh{Practice solutions}

\textbf{Level 0:} It is a scatter plot. Width is on the x-axis and area is on
the y-axis.

\textbf{Level 2}

\begin{lstlisting}[
    language=python,
    style=block,
    caption={One solution to Level 2},
    label={c13_solution_level2},
    captionpos=b
]
import matplotlib.pyplot as plt

hours = [8, 10, 12, 14, 16, 18]
temperatures_c = [19.0, 20.5, 22.0, 23.0, 22.0, 20.5]

fig, ax = plt.subplots(layout="constrained")
ax.plot(hours, temperatures_c, color="tab:blue", marker="o")
ax.set(
    title="Temperature through the day",
    xlabel="Hour",
    ylabel="Temperature (degrees C)",
)
ax.grid(axis="y", alpha=0.25)

plt.show()
plt.close(fig)
\end{lstlisting}

\textbf{Level 3}

\begin{lstlisting}[
    language=python,
    style=block,
    caption={One solution to Level 3},
    label={c13_solution_level3},
    captionpos=b
]
import matplotlib.pyplot as plt
import pandas as pd

rooms = pd.DataFrame({
    "room": ["A", "B", "C", "D", "E"],
    "area_m2": [18.0, 24.0, 31.0, 26.0, 21.0],
})

fig, axes = plt.subplots(1, 2, figsize=(9, 4), layout="constrained")
axes[0].bar(rooms["room"], rooms["area_m2"], color="tab:blue")
axes[0].set(
    title="Area by room",
    xlabel="Room",
    ylabel="Area (m2)",
)

axes[1].hist(rooms["area_m2"], bins=4, color="tab:orange", edgecolor="black")
axes[1].set(
    title="Area distribution",
    xlabel="Area (m2)",
    ylabel="Number of rooms",
)

for ax in axes:
    ax.grid(axis="y", alpha=0.25)
    ax.set_axisbelow(True)

plt.show()
plt.close(fig)
\end{lstlisting}

\textbf{Level 4}

\begin{lstlisting}[
    language=python,
    style=block,
    caption={One solution to Level 4},
    label={c13_solution_level4},
    captionpos=b
]
import matplotlib.pyplot as plt

def plot_two_series(x, first, second, labels):
    fig, ax = plt.subplots(layout="constrained")
    ax.plot(
        x,
        first,
        color="tab:blue",
        marker="o",
        linestyle="-",
        label=labels[0],
    )
    ax.plot(
        x,
        second,
        color="tab:orange",
        marker="s",
        linestyle="--",
        label=labels[1],
    )
    ax.set(
        title="Temperature comparison",
        xlabel="Hour",
        ylabel="Temperature (degrees C)",
    )
    ax.legend()
    ax.grid(axis="y", alpha=0.25)
    return fig

hours = [8, 10, 12, 14]
north_c = [19.0, 20.0, 21.0, 21.5]
south_c = [19.5, 21.0, 23.0, 24.0]
figure = plot_two_series(
    hours,
    north_c,
    south_c,
    labels=("North room", "South room"),
)
plt.close(figure)
\end{lstlisting}

\hh{Chapter checkpoint}

\begin{enumerate}
    \item What is the difference between a Figure and an Axes?
    \item Which plot suits change through an ordered sequence?
    \item Which plot suits comparison among separate categories?
    \item Which plot shows a relationship between paired numerical variables?
    \item What does a histogram bin represent?
    \item Why should a bar chart usually begin at zero?
    \item Why add markers or line styles when colors already differ?
    \item What does \c{layout="constrained"} help prevent?
    \item Which object should call \c{savefig()} in this chapter's pattern?
    \item What should you do after saving a figure?
\end{enumerate}

\hh{Checkpoint answers}

\begin{enumerate}
    \item The Figure is the complete output; an Axes is one plotting region within it.
    \item A line plot.
    \item A bar chart.
    \item A scatter plot.
    \item An interval used to count numerical observations.
    \item Bar length encodes value, so a nonzero baseline can exaggerate differences.
    \item They preserve distinctions when color is difficult to perceive or unavailable.
    \item Overlap among titles, labels, and plotting regions.
    \item The Figure object.
    \item Open and inspect the saved file for labels, cropping, and readability.
\end{enumerate}

\begin{tcolorbox}[title=Part IV summary,colframe=orange,colback=white,arc=2mm]
NumPy organizes numerical arrays, pandas adds labels and table operations, and
Matplotlib communicates selected patterns. Begin with a question, prepare the
data visibly, choose a plot whose structure matches the variables, and label
the result so another reader can understand it without seeing the code.
\end{tcolorbox}

